\section{Preliminaries}

\subsection{Multi-turn Tool-Integrated Reasoning}

We consider the post-training of a small language model (SLM) for multi-turn tool-integrated reasoning (TIR). Given an input $x$, the model interacts with an external environment over multiple reasoning steps. At each step $k$, the model generates a response $y_k$, which may include natural-language reasoning, a tool invocation, or a final answer. If a tool is invoked, the environment returns an observation $o_k$, which is appended to the context and conditions subsequent generations.

A trajectory is defined as:
\begin{equation}
    \tau = (x, y_1, o_1, \ldots, y_K, o_K, y_{K+1}),
\end{equation}
where $y_{K+1}$ denotes the final response. The policy $\pi_\theta$ generates only model tokens, while observations $\{o_k\}$ are provided by the environment. Let $y_t$ denote a generated token and $y_{<t}$ its prefix, which may include both model outputs and tool observations.


\subsection{Group Relative Policy Optimization}

Group Relative Policy Optimization (GRPO)~\citep{grpo} is a reinforcement learning algorithm that updates the policy using relative rewards within a group of sampled trajectories. We assume access to an outcome-level reward function $r(\tau)$ defined on complete trajectories. For each input $x$, a group of trajectories $\{\tau_i\}_{i=1}^G$ is sampled from the old policy $\pi_{\theta_{\mathrm{old}}}$, each receiving reward $r_i = r(\tau_i)$.

The group-relative advantage is computed as:
\begin{equation}
    \hat{A}_i =
    \frac{r_i - \mathrm{mean}(\{r_j\})}
    {\mathrm{std}(\{r_j\}) + \epsilon_A}.
\label{eq:advantage}
\end{equation}

Let $\rho_{i,t}(\theta) = \frac{\pi_\theta(y_{i,t} \mid y_{i,<t})}{\pi_{\theta_{\mathrm{old}}}(y_{i,t} \mid y_{i,<t})}$ denote the token-level importance ratio. The objective is defined as:
\begin{equation}
\mathcal{L}_{\mathrm{GRPO}}
=
-
\mathbb{E}
\left[
\frac{1}{G}
\sum_{i=1}^{G}
\frac{1}{|\mathcal{T}_i|}
\sum_{t \in \mathcal{T}_i}
\min
\left(
\rho_{i,t}(\theta)\hat{A}_i,
\mathrm{clip}(\rho_{i,t}(\theta), 1-\epsilon, 1+\epsilon)\hat{A}_i
\right)
\right],
\label{eq:grpo}
\end{equation}
where $\mathcal{T}_i$ denotes model-generated token positions.
This objective provides an on-policy learning signal based on relative trajectory performance, but supplies only sparse outcome-level rewards.

\subsection{On-policy Distillation}

On-policy distillation (OPD)~\citep{opd,lu2025onpolicydistillation} is a post-training paradigm that provides dense token-level supervision on student-generated trajectories by aligning the student policy with a teacher distribution. Given a trajectory $y \sim \pi_\theta$, the OPD objective is defined as:
\begin{equation}
\mathcal{L}_{\mathrm{OPD}}^{\mathrm{}}
=
\mathbb{E}
\left[
\sum_{t \in \mathcal{T}_i}
\left(
\log \pi_\theta(y_t \mid y_{<t})
-
\log \pi_{\mathrm{teacher}}(y_t \mid y_{<t})
\right)
\right],
\end{equation}

where $\pi_\theta$ is the student policy, $\pi_{\mathrm{teacher}}$ is the teacher model, and $\mathcal{T}_i$ denotes the set of model-generated token positions in the sampled trajectory. This objective corresponds to a sampled estimator of the reverse KL divergence between the student and teacher policies on student-visited states. 







